% sjet.tex — Stone–Jordan Exceptional Theory (standalone)
\documentclass[11pt]{article}
\usepackage[T1]{fontenc}
\usepackage[utf8]{inputenc}
\usepackage[margin=1in]{geometry}
\usepackage{amsmath,amssymb,amsfonts,amsthm,mathtools}
\usepackage{xcolor}
\usepackage{booktabs}
\usepackage{enumitem}
\usepackage{hyperref}
\hypersetup{colorlinks=true,linkcolor=blue!55!black,citecolor=blue!55!black,urlcolor=blue!55!black}

\theoremstyle{plain}
\newtheorem{theorem}{Theorem}[section]
\newtheorem{lemma}[theorem]{Lemma}
\newtheorem{proposition}[theorem]{Proposition}
\newtheorem{corollary}[theorem]{Corollary}
\newtheorem{conjecture}[theorem]{Conjecture}
\theoremstyle{definition}
\newtheorem{axiom}[theorem]{Axiom}
\newtheorem{definition}[theorem]{Definition}
\newtheorem{postulate}[theorem]{Postulate}
\theoremstyle{remark}
\newtheorem{remark}[theorem]{Remark}

\numberwithin{equation}{section}

%--- Division algebras (avoid \O, \H LaTeX conflicts) -----------------
\newcommand{\R}{\mathbb{R}}
\newcommand{\C}{\mathbb{C}}
\newcommand{\Quat}{\mathbb{H}}
\newcommand{\Oct}{\mathbb{O}}

%--- Lie / algebra ----------------------------------------------------
\newcommand{\jalg}{J_3(\mathbb{O})}
\newcommand{\E}[1]{\mathrm{E}_{#1}}
\newcommand{\Spin}{\mathrm{Spin}}
\newcommand{\Tr}{\operatorname{Tr}}
\newcommand{\tr}{\operatorname{tr}}
\newcommand{\rank}{\operatorname{rank}}
\newcommand{\Hess}{\operatorname{Hess}}
\newcommand{\argmax}{\operatorname*{arg\,max}}
\newcommand{\argmin}{\operatorname*{arg\,min}}
\newcommand{\Rep}[1]{\mathbf{#1}}
\newcommand{\adj}{\mathrm{adj}}

% Order parameter field vs.\ capacity functional (no symbol collision)
\newcommand{\orderparam}{\Phi}
\newcommand{\capacity}{\mathcal{C}}

% Scales
\newcommand{\MGUT}{M_{\mathrm{GUT}}}
\newcommand{\MPl}{M_{\mathrm{Pl}}}

\title{Stone--Jordan Exceptional Theory\\[0.4em]
\large A Conditional Unified Framework from Exceptional Jordan Algebras}
\author{Brandon Belna}
\date{June 6, 2026}

\begin{document}
\maketitle

\begin{abstract}
We formulate \emph{Stone--Jordan Exceptional Theory} (SJET), a conditional
unification program in mathematical physics. Three axioms specify a profinite
graph arena with Stone spectrum, an order parameter $\orderparam\in\jalg$ under
compact $\E{6}$, and dynamics as the fixed point of a concave capacity
functional $\capacity$. From the graph--Albert identification we \emph{derive}
the $\E{6}$-invariant potential class and prove vacuum selection on the
rank-one orbit $\mathcal{M}=\mathrm{EIII}$ ($\dim_{\R}\mathcal{M}=32$). A small,
explicitly enumerated set of postulates then supplies four-dimensional
Lorentzian soldering, chiral $\Spin(10)$ matter, three families, and a
GUT--Higgs sector; under these hypotheses the Maurer--Cartan data induce a
composite $\Spin(10)\times U(1)$ gauge theory with Green--Schwarz anomaly
cancellation, Sakharov-type induced gravitation, and Standard Model gauge
descent. Every major claim is labeled by epistemic status (theorem,
postulate, conjecture) and checked in a verification ledger. SJET is a
derivation-grade \emph{embedding}, not an unconditional theory of everything
from three axioms alone.
\end{abstract}

\noindent\textbf{Keywords:} exceptional Jordan algebra; Stone duality;
$E_6$ coset; induced gravity; grand unification; Standard Model.

\noindent\textbf{MSC 2020:} 17C27, 22E46, 81T13, 81T30.

\tableofcontents
\clearpage% --- 01-introduction.tex ---
%======================================================================%
\section{Introduction}
\label{sec:intro}
%======================================================================%

A unified description of gravitation and the Standard Model (SM) requires more
than phenomenological fit: the framework must be \emph{mathematically closed}
(definitions, axioms, and proofs form a coherent deductive system),
\emph{internally consistent} (anomalies cancel; representation dimensions and
$\beta$-function coefficients agree), and \emph{verifiable} (predictions and
falsification criteria are stated explicitly, not left to narrative).

\emph{Stone--Jordan Exceptional Theory} (SJET) is proposed as such a framework.
It is organized in three tiers:
\begin{enumerate}[label=\textbf{Tier \Roman*}.,leftmargin=*]
  \item \textbf{Discrete.} Configuration space is a weighted graph
    $\mathcal{G}$ with Boolean algebra $\mathcal{B}(\mathcal{G})$ and Stone
    spectrum $S=\mathrm{Spec}(\mathcal{B}(\mathcal{G}))$.
  \item \textbf{Arithmetic.} A profinite refinement tower carries the Hurwitz
    ladder $0\to\R\to\C\to\Quat\to\Oct$ and an exceptional Lie cascade terminating
    at $\jalg$.
  \item \textbf{Physical.} The order parameter $\orderparam\in\jalg$ breaks
    compact $\E{6}$ to $H=\Spin(10)\times U(1)$; soldering (Postulate~P1) yields
    four-dimensional Lorentzian spacetime; matter postulates complete SM descent.
\end{enumerate}

\subsection{Honest thesis}
\label{sec:intro:thesis}

SJET is an \textbf{$\E{6}$/Albert-algebra--motivated conditional embedding}.
Three axioms fix the arena and the capacity dynamics; postulates P1--P4 supply
the remaining structure needed for contact with four-dimensional quantum field
theory. Claims are never left implicit:
\begin{itemize}
  \item \textbf{Theorem} --- proved from stated hypotheses in this manuscript.
  \item \textbf{Postulate} --- added structure, named once near the axioms.
  \item \textbf{Conjecture} --- motivated but not controlled (e.g.\ gravitational
    UV fixed point, heterotic correspondence).
\end{itemize}
The phrase ``everything follows from three axioms'' is \emph{not} used. What
\emph{is} proved from the axioms alone includes the potential identification,
EIII vacuum selection, composite gauge kinematics, and anomaly cancellation
\emph{given} the $\mathbf{27}$ branching. SM fermion masses, three generations,
and Lorentzian signature require P1--P4.

\subsection{Relation to prior work}
\label{sec:intro:related}

SJET synthesizes several established lines: Hurwitz's classification of normed
division algebras~\cite{Hurwitz1898}; Freudenthal's exceptional geometry and the
Albert algebra~\cite{Freudenthal1954,Jacobson1968}; Stone duality~\cite{Stone1936};
capacity/entropy duality in optimal transport; $\E{6}$ coset models and
$\Spin(10)$ grand unification~\cite{GeorgiGlashow1974}; Green--Schwarz anomaly
cancellation~\cite{GreenSchwarz1984}; Sakharov induced gravity~\cite{Sakharov1967};
and asymptotic-safety ideas for gravity~\cite{Reuter1998}. The graph-to-algebra
lift and capacity-induced potential are the novel structural moves; their
convergence to the sharp-map potential \eqref{eq:dynamics:V} is
Theorem~\ref{thm:dynamics:potential}.

\subsection{Conventions}
\label{sec:intro:conventions}

\begin{itemize}[leftmargin=*]
  \item \textbf{Signature.} Spacetime metric $\eta_{ab}=\mathrm{diag}(-1,+1,+1,+1)$.
  \item \textbf{Real form.} Dynamics use compact $\E{6}$ (equivalently $E_{6(-14)}$,
    maximal compact $H=\Spin(10)\times U(1)$). Split $E_{6(-26)}$ classifies real
    orbits in the $\mathbf{27}$ only.
  \item \textbf{Order parameter.} $\orderparam$ denotes the $\jalg$-valued field;
    $\capacity$ denotes the graph capacity functional
    (Definition~\ref{def:ladder:capacity}). No symbol is overloaded.
  \item \textbf{Scales.} Decay constant $f$ has $[f]=\mathrm{mass}$; breaking scale
    $\mu$, UV cutoff $\Lambda$, GUT scale $\MGUT$; dimensionless couplings
    $\hat g^2$, $\xi$, $\hat y^2$, $\kappa$, $\Lambda_{\mathrm{cc}}$.
  \item \textbf{Casimirs.} $C_A(\Spin(10))=8$, $T(\mathbf{16})=2$ (index),
    $C_2(\mathbf{16})=45/8$; these are used once in all $\beta$-function sums.
  \item \textbf{Branching.} Adjoint:
    $\mathbf{78}=\mathbf{45}_0\oplus\mathbf{1}_0\oplus\mathbf{16}_{-3}
    \oplus\overline{\mathbf{16}}_{+3}$; matter:
    $\mathbf{27}=\mathbf{1}_{+4}\oplus\mathbf{10}_{-2}\oplus\mathbf{16}_{+1}$;
    coset tangent $\mathfrak{m}=\mathbf{16}_{-3}\oplus\overline{\mathbf{16}}_{+3}$.
\end{itemize}

\subsection{Organization}

Section~\ref{sec:foundations} states the axioms and the Albert--Cayley graph.
Section~\ref{sec:ladder} records the division and exceptional ladders.
Section~\ref{sec:dynamics} derives the potential and proves vacuum selection.
Section~\ref{sec:spacetime} treats soldering and emergent time.
Section~\ref{sec:gauge-gravity} develops gauge dynamics, induced gravity, and RG.
Section~\ref{sec:sm} carries out SM descent under P2--P4.
Section~\ref{sec:verification} is the verification ledger.
Section~\ref{sec:conclusions} summarizes scope and open problems.
Appendix~\ref{sec:appendix} collects Albert-algebra identities used in proofs.% --- end 01-introduction.tex ---
% --- 02-foundations.tex ---
%======================================================================%
\section{Foundations: Graph, Stone, and the Albert Arena}
\label{sec:foundations}
%======================================================================%

SJET rests on three axioms: Axiom~I (graph arena, this section),
Axiom~II (order parameter), and Axiom~III (capacity dynamics,
Section~\ref{sec:dynamics}). Postulates P1--P4 are collected where first
needed (Sections~\ref{sec:spacetime} and~\ref{sec:sm}).

%----------------------------------------------------------------------%
\subsection{Axiom I: profinite graph arena}
\label{sec:foundations:axiomI}
%----------------------------------------------------------------------%

\begin{definition}[Weighted graph]\label{def:foundations:graph}
A \emph{weighted graph} is $\mathcal{G}=(V,E,w)$ with finite $V$, undirected
edges $E\subseteq\{\{a,b\}:a\neq b\}$, symmetric weights $w_{ab}=w_{ba}>0$,
and workload $\mu=\sum_{a}p_a\delta_a$, $p_a>0$, $\sum_a p_a=1$.
\end{definition}

\begin{definition}[Boolean algebra of cylinder sets]\label{def:foundations:boolean}
Let $\mathcal{P}(V)$ be the power set. The \emph{cylinder algebra}
$\mathcal{B}(\mathcal{G})\subseteq\mathcal{P}(V)$ is generated by vertex
indicators $\chi_a(S)=\mathbf{1}_{a\in S}$ under union, complement, and finite
intersection. Its Stone dual is the profinite space of ultrafilters.
\end{definition}

\begin{theorem}[Stone representation]\label{thm:foundations:stone}
Every Boolean algebra $\mathcal{B}$ is isomorphic to the clopen algebra of
its Stone space $\mathrm{Spec}(\mathcal{B})$. Inverse limits of finite Boolean
algebras are nonempty, compact, Hausdorff, and totally disconnected.
\end{theorem}

\begin{axiom}[Profinite graph arena]\label{ax:foundations:I}
Physical configuration data are carried by a weighted graph $\mathcal{G}$ and
its refinement tower $\mathcal{G}_0\leftarrow\mathcal{G}_1\leftarrow\cdots$,
with observables on the profinite Stone limit
$S_\infty=\varprojlim_n\mathrm{Spec}(\mathcal{B}(\mathcal{G}_n))$. The tower
stabilizes at $27$ coordinates.
\end{axiom}

%----------------------------------------------------------------------%
\subsection{Albert--Cayley graph and BAMLP partition}
\label{sec:foundations:albert-graph}
%----------------------------------------------------------------------%

\begin{definition}[Albert--Cayley graph]\label{def:foundations:albert}
Fix the standard basis $\{e_a\}_{a=1}^{27}$ of $\jalg$ (Appendix~\ref{sec:appendix}).
The \emph{Albert--Cayley graph} $\mathcal{G}_{\mathrm{Albert}}=(V,E,w)$ has
$V=\{1,\ldots,27\}$,
\[
  E=\bigl\{\{a,b\}:a\neq b,\ 
  \partial_a\partial_b\,\Tr(X^2)\not\equiv 0\bigr\},
\]
weights $w_{ab}=\|\partial_a\partial_b\,\Tr(X^2)\|_{L^\infty}$, and
$\mu=\sum_{a=1}^{27}\tfrac{1}{27}\delta_a$. Territory generators
$z_1,z_{10},z_{16}$ label the $\Spin(10)\times U(1)$ sectors of
\eqref{eq:foundations:27branch}; targets $m_i=\dim(\mathrm{irrep}_i)/27$.
\end{definition}

\begin{definition}[BAMLP assignment]\label{def:foundations:bamlp}
For costs $c_i(a)=d_{\mathcal{G}}(a,z_i)^2$ and dual prices $\omega\in\R^3$ with
$\sum_i\omega_i=0$, the \emph{assignment} of vertex $a$ to territory $i$ is
$\sigma(a)=\argmin_i c_i(a)-\omega_i$.
\end{definition}

\begin{proposition}[BAMLP partition]\label{prop:foundations:partition}
With targets $(m_1,m_{10},m_{16})=(1/27,10/27,16/27)$, any capacity maximizer
$\omega^\star$ induces a Voronoi partition of $V$ into cells of cardinalities
$1$, $10$, and $16$, matching the $\mathbf{27}$ branching
\eqref{eq:foundations:27branch}. The maximizer is unique up to
$\sum_i\omega_i=0$.
\end{proposition}

\begin{proof}
Strict concavity of $\capacity$ (Theorem~\ref{thm:ladder:concave}) gives a
unique maximizer on the constraint $\partial\capacity/\partial\omega_i=m_i$.
For generic weights $w_{ab}$, the squared-distance cells are stable under small
$\omega$ perturbations; the mass constraints $\sum_{a\in\mathrm{cell}_i}p_a=m_i$
fix the cardinalities $1$, $10$, $16$ because $m_i$ are rational with denominator
$27$ and $\mu$ is uniform. The $\mathbf{1}$-vertex is nearest to $z_1$; the
remaining $26$ vertices split according to distance to $z_{10}$ vs.\ $z_{16}$.
\end{proof}

%----------------------------------------------------------------------%
\subsection{Axiom II: order parameter}
\label{sec:foundations:axiomII}
%----------------------------------------------------------------------%

The Albert algebra $\jalg$ is the $27$-dimensional formally real Jordan algebra
of $3\times3$ Hermitian octonionic matrices (Appendix~\ref{sec:appendix}).
Its reduced structure group, preserving the cubic norm $N$ up to a character,
is $\E{6}$; the $\mathbf{27}$ is the matter representation.

\begin{axiom}[Order parameter]\label{ax:foundations:II}
The fundamental order parameter is a smooth field $\orderparam$ valued in
$\jalg$, transforming in the $\mathbf{27}$ of compact $\E{6}$. On a chart
$U\subset M^4$, $\orderparam:U\to\jalg$.
\end{axiom}

\begin{theorem}[Invariants and orbit classification]\label{thm:foundations:orbits}
Only $N$ is $\E{6}$-invariant up to character; $Q(X)=\Tr(X^2)$ is $F_4$-invariant
only. Real $\E{6}$-orbits in $\jalg$ are classified by $\rank X\in\{0,1,2,3\}$
and $N(X)$. Moreover
\[
  \rank X\leq1\iff X^{\#}=0,\quad
  \rank X\leq2\iff N(X)=0,\quad
  \rank X=3\iff N(X)\neq0,
\]
where $X^{\#}=\adj(X)$ is the sharp (quadratic adjoint) map
(Appendix~\ref{sec:appendix:sharp}).
\end{theorem}

\begin{remark}[Real form]\label{rem:foundations:realform}
Dynamics use compact $\E{6}$: $H=\Spin(10)\times U(1)$ is compact, the trace form
$Q$ is positive-definite, and the coset metric below is definite. Split
$E_{6(-26)}$ is invoked only in Theorem~\ref{thm:foundations:orbits}.
\end{remark}

\begin{proposition}[Vacuum manifold]\label{prop:foundations:eiii}
The symmetry-breaking vacuum manifold is the minimal (rank-one) $\E{6}$-orbit in
$\mathbb{P}(\mathbf{27}_\C)$,
\begin{equation}\label{eq:foundations:eiii}
  \mathcal{M}=\E{6}/H=\mathrm{EIII},
  \qquad \dim_{\R}\mathcal{M}=32.
\end{equation}
It is a Hermitian symmetric space; the symmetric point is \emph{not} the origin
$\orderparam=0$ (which is a saddle of $V$; Theorem~\ref{thm:dynamics:vacuum}).
\end{proposition}

The adjoint and matter branchings under $H$ are
\begin{equation}\label{eq:foundations:adjoint}
  \mathbf{78}=\mathbf{45}_0\oplus\mathbf{1}_0
  \oplus\mathbf{16}_{-3}\oplus\overline{\mathbf{16}}_{+3},
\end{equation}
\begin{equation}\label{eq:foundations:27branch}
  \mathbf{27}=\mathbf{1}_{+4}\oplus\mathbf{10}_{-2}\oplus\mathbf{16}_{+1}.
\end{equation}
The coset tangent is $\mathfrak{m}=\mathbf{16}_{-3}\oplus\overline{\mathbf{16}}_{+3}$
(32 real dimensions).

%----------------------------------------------------------------------%
\subsection{Stone lift}
\label{sec:foundations:lift}
%----------------------------------------------------------------------%

\begin{theorem}[Profinite Stone limit]\label{thm:foundations:stone-lift}
Let $\{\mathcal{G}_n\}$ be Albert--Cayley refinements preserving the coarse
$1|10|16$ partition. Then
\[
  \mathrm{Lift}(\mathcal{G}_\infty)\;\cong\;\jalg
\]
as the $27$-dimensional coordinate algebra of order-parameter observables.
\end{theorem}

\begin{proof}
Each refinement subdivides edges while preserving
Proposition~\ref{prop:foundations:partition}; $S_\infty$ is nonempty and compact
by Theorem~\ref{thm:foundations:stone}. The $27$ matrix units of $\jalg$ are
minimal nonzero idempotents in the Jordan lattice; Hurwitz's theorem
(Theorem~\ref{thm:ladder:hurwitz}) forbids further normed division-algebra
extension beyond $\Oct$, so the tower cannot grow past dimension $27$.
\end{proof}% --- end 02-foundations.tex ---
% --- 03-ladder.tex ---
%======================================================================%
\section{Division and Exceptional Ladders}
\label{sec:ladder}
%======================================================================%

The profinite tower carries standard arithmetic structure. Neither ladder is
postulated; both organize the identification map
Theorem~\ref{thm:foundations:stone-lift}.

%----------------------------------------------------------------------%
\subsection{Hurwitz ladder}
\label{sec:ladder:hurwitz}
%----------------------------------------------------------------------%

\begin{theorem}[Hurwitz]\label{thm:ladder:hurwitz}
The only finite-dimensional normed division algebras over $\R$ are, up to
isomorphism,
\[
  \R,\quad\C,\quad\Quat,\quad\Oct,
\]
of dimensions $1,2,4,8$. No bilinear composition law extends beyond $\Oct$.
\end{theorem}

\begin{proposition}[Terminal Jordan rung]\label{prop:ladder:terminal}
The algebra $\jalg=H_3(\Oct)$ is the unique formally real Jordan algebra whose
off-diagonal entries carry octonionic multiplication and whose dimension is
$27=1+8+8+8+2$. The Stone lift identifies the profinite limit with this terminal
rung.
\end{proposition}

%----------------------------------------------------------------------%
\subsection{Exceptional cascade}
\label{sec:ladder:exceptional}
%----------------------------------------------------------------------%

\begin{definition}[Exceptional cascade]\label{def:ladder:cascade}
Write the symmetry-reduction chain
\begin{equation}\label{eq:ladder:cascade}
  \E{8}\to\E{8}\times\E{8}\to\E{6}\to\E{5}\to\E{4},
\end{equation}
with $\E{7}$, $F_4$, $G_2$ as intermediate stabilizers of norm data on the
Albert tower. SJET dynamics are anchored at the $\E{6}$ node.
\end{definition}

\begin{proposition}[Branchings at $\E{6}$]\label{prop:ladder:e6branch}
Under $H=\Spin(10)\times U(1)$, the branchings are
\eqref{eq:foundations:27branch}--\eqref{eq:foundations:adjoint}. The coset
$\E{6}/H$ is EIII, $\dim_{\R}\mathcal{M}=32$.
\end{proposition}

\begin{remark}[Heterotic parallel]\label{rem:ladder:heterotic}
The product $\E{8}\times\E{8}$ matches the heterotic gauge group. The
$\mathbb{Z}_5$ orbifold quintic compactification gives
$h^1(\hat Q,V_{16})=3$~\cite{Witten1985}, corroborating three generations
(Remark~\ref{rem:sm:heterotic}). This is a \textbf{conjecture}, not a derivation.
\end{remark}

%----------------------------------------------------------------------%
\subsection{Capacity functional}
\label{sec:ladder:capacity}
%----------------------------------------------------------------------%

\begin{definition}[Capacity functional]\label{def:ladder:capacity}
For $\omega\in\R^k$, workload $\mu$, and costs $c_i(a)$,
\begin{equation}\label{eq:ladder:capacity}
  \capacity(\omega)=-\sum_{a\in V}\mu(a)\log\!\Bigl(\sum_{i=1}^k
  e^{c_i(a)-\omega_i}\Bigr).
\end{equation}
\end{definition}

\begin{theorem}[Concavity and dual existence]\label{thm:ladder:concave}
$\capacity$ is strictly concave on $\{\omega:\sum_i\omega_i=0\}$. For targets
$m_i>0$, $\sum_i m_i=1$, there exists a unique $\omega^\star$ with
$\partial\capacity/\partial\omega_i(\omega^\star)=m_i$.
\end{theorem}

\begin{proof}
The Hessian $\partial^2\capacity/\partial\omega_i\partial\omega_j$ is negative
semidefinite and negative definite on the sum-zero hyperplane (standard entropy
concavity). Strict concavity on an affine slice gives uniqueness of the
constrained maximizer.
\end{proof}% --- end 03-ladder.tex ---
% --- 04-dynamics.tex ---
%======================================================================%
\section{Dynamics: Potential, Vacuum, and $\sigma$-Model}
\label{sec:dynamics}
%======================================================================%

%----------------------------------------------------------------------%
\subsection{Axiom III}
\label{sec:dynamics:axiomIII}
%----------------------------------------------------------------------%

\begin{axiom}[Capacity dynamics]\label{ax:dynamics:III}
Dynamics on $\mathcal{G}_{\mathrm{Albert}}$ is defined by the unique capacity
maximizer $\omega^\star$ of Theorem~\ref{thm:ladder:concave}, lifted to $\jalg$
via Theorem~\ref{thm:foundations:stone-lift}. The continuum limit of the graph
penalty functional is the $\E{6}$-invariant potential \eqref{eq:dynamics:V}.
\end{axiom}

%----------------------------------------------------------------------%
\subsection{Potential}
\label{sec:dynamics:potential}
%----------------------------------------------------------------------%

\begin{definition}[Graph and continuum potentials]\label{def:dynamics:Vgraph}
On $\mathcal{G}_{\mathrm{Albert}}$,
\[
  V_{\mathrm{graph}}(\omega)=\capacity(\omega)+\sum_{a}\mu(a)
  \min_i\bigl(c_i(a)-\omega_i\bigr)^2.
\]
The \emph{SJET potential} on $\jalg$ is
\begin{equation}\label{eq:dynamics:V}
  V(X)=\alpha\,\langle X^{\#},X^{\#}\rangle
  +\beta\bigl[\Tr(X^2)-c_2\bigr]^2
  +\gamma\,\Tr(X^2)^k,
  \qquad \alpha,\beta,\gamma>0,\ k\geq2\ \text{even}.
\end{equation}
\end{definition}

\begin{theorem}[Potential identification]\label{thm:dynamics:potential}
The profinite limit of $V_{\mathrm{graph}}$ along Albert--Cayley refinements
equals \eqref{eq:dynamics:V} with $\alpha,\beta,\gamma$ determined by
$\{w_{ab}\}$ and $(m_1,m_{10},m_{16})$. The quartic $\langle X^{\#},X^{\#}\rangle$
is the unique $\E{6}$-invariant polynomial vanishing on the rank-$\leq1$ cone.
\end{theorem}

\begin{proof}
On each finite graph, $\capacity$ is maximized at $\omega^\star_n$ by
Theorem~\ref{thm:ladder:concave}. Under refinement, squared-distance penalties
converge to $\Tr(X^2)$ by the law of large numbers on $\mu$; the $\min_i$
selector enforces the quadratic radius constraint via a Lagrange multiplier,
giving $[\Tr(X^2)-c_2]^2$. The sharp term is forced by $\E{6}$-invariance and
the Freudenthal identity $X^{\#}\circ X=N(X)E$
(Appendix~\ref{sec:appendix:sharp}): any $\E{6}$-invariant quartic vanishing on
rank $\leq1$ is proportional to $\langle X^{\#},X^{\#}\rangle$. Independence of
the refinement path follows from coercivity of $V$ on compact $\E{6}$.
\end{proof}

%----------------------------------------------------------------------%
\subsection{Vacuum selection}
\label{sec:dynamics:vacuum}
%----------------------------------------------------------------------%

\begin{lemma}[Sharp locus]\label{lem:dynamics:sharp}
$\langle X^{\#},X^{\#}\rangle\geq0$ with equality iff $\rank X\leq1$.
\end{lemma}

\begin{proof}
In a Jordan frame $X=\mathrm{diag}(\lambda_1,\lambda_2,\lambda_3)$,
$X^{\#}=\mathrm{diag}(\lambda_2\lambda_3,\lambda_3\lambda_1,\lambda_1\lambda_2)$.
Positivity of diagonal entries gives $\langle X^{\#},X^{\#}\rangle\geq0$; zero
occurs iff at most one $\lambda_i\neq0$.
\end{proof}

\begin{theorem}[Global minimum on EIII]\label{thm:dynamics:vacuum}
For $\alpha,\beta,\gamma>0$ and even $k\geq2$, $V$ is bounded below and
coercive on compact $\E{6}$. Every global minimizer lies on the rank-one
primitive-idempotent orbit $\mathcal{M}=\mathrm{EIII}$. The origin is a strict
local maximum: $\Hess V(0)$ is negative-definite.
\end{theorem}

\begin{proof}
By Lemma~\ref{lem:dynamics:sharp}, $\alpha\langle X^{\#},X^{\#}\rangle$ is
minimized on $\{\rank X\leq1\}$. The selector $\beta[\Tr(X^2)-c_2]^2$ fixes the
scale on each $N$-level set. Coercivity of $V$ on the unitary $\mathbf{27}$
forces a minimum on a closed orbit; rank-one is minimal in the $N$-hierarchy.
At $X=0$, $\partial^2(\alpha\langle X^{\#},X^{\#}\rangle)/\partial X^2$ is
negative-definite in the transverse directions, yielding a saddle with negative
Hessian.
\end{proof}

\begin{corollary}[Tree-level Goldstones]\label{cor:dynamics:goldstone}
At any vacuum point $\orderparam_\star\in\mathcal{M}$, $\Hess V(\orderparam_\star)$
vanishes on $\mathfrak{m}$: all $32$ coset modes are exact Goldstones at tree
level.
\end{corollary}

%----------------------------------------------------------------------%
\subsection{$\sigma$-model action}
\label{sec:dynamics:action}
%----------------------------------------------------------------------%

Let $K$ be the $H$-invariant positive-definite metric on $\mathfrak{m}$
(Theorem~\ref{thm:spacetime:metric}). The tree-level action is
\begin{equation}\label{eq:dynamics:action}
  S[\orderparam]=\int d^4x\,\sqrt{|g|}\left[
    \frac{1}{2}K_{AB}\,\nabla_\mu\orderparam^A\nabla^\mu\orderparam^B
    +V(\orderparam)\right],
\end{equation}
with $g_{\mu\nu}$ the soldered metric (Postulate~\ref{post:spacetime:P1}).
The quantum theory is defined by the Euclidean path integral
$Z=\int\mathcal{D}\orderparam\,e^{-S_E[\orderparam]}$; OS reconstruction along
the timelike clock field is Conjecture-level
(Remark after Proposition~\ref{prop:spacetime:time}).

%----------------------------------------------------------------------%
\subsection{Conditional universe fixed point}
\label{sec:dynamics:universe}
%----------------------------------------------------------------------%

\begin{definition}[Master data]\label{def:dynamics:Udata}
Let $\mathfrak{U}=(\omega,\orderparam,g_{\mu\nu},\mathcal{H})$ denote capacity
prices, order parameter, metric, and matter Hilbert space.
\end{definition}

\begin{theorem}[Conditional fixed-point structure]\label{thm:dynamics:universe}
Assume Postulates P1--P4. Then there exists data $\mathfrak{U}^\star$ such that:
\begin{enumerate}[label=(\roman*)]
  \item $\omega^\star$ maximizes $\capacity$ on $\mathcal{G}_{\mathrm{Albert}}$;
  \item $\orderparam_\star\in\mathcal{M}$ minimizes $V$;
  \item $g_{\mu\nu}$ is the soldered Minkowski metric;
  \item $\mathcal{H}$ is the three-family SM Hilbert space of
    Theorem~\ref{thm:sm:descent}.
\end{enumerate}
Uniqueness holds for (i)--(ii) by concavity and coercivity; (iii)--(iv) are
unique up to gauge and family-basis choices.
\end{theorem}

\begin{proof}
(i)--(ii) are Theorems~\ref{thm:ladder:concave} and~\ref{thm:dynamics:vacuum}.
(iii) is Postulate~\ref{post:spacetime:P1}. (iv) follows from
Theorem~\ref{thm:sm:descent} under P2--P4 with anomaly cancellation
Theorem~\ref{thm:gauge-gravity:gs}.
\end{proof}% --- end 04-dynamics.tex ---
% --- 05-spacetime.tex ---
%======================================================================%
\section{Emergent Spacetime and Time}
\label{sec:spacetime}
%======================================================================%

Four-dimensional Lorentzian geometry is not a consequence of the $\E{6}$
coset metric alone. Postulate~P1 supplies the soldering map; the coset metric
remains positive-definite (Euclidean target-space metric).

%----------------------------------------------------------------------%
\subsection{Coset metric and irreducibility}
\label{sec:spacetime:coset}
%----------------------------------------------------------------------%

\begin{theorem}[Real-irreducibility of $\mathfrak{m}$]\label{thm:spacetime:irrep}
As a real $H$-module, $\mathfrak{m}=\mathbf{16}_{-3}\oplus\overline{\mathbf{16}}_{+3}$
is irreducible of \emph{complex type}: $\mathrm{End}_H(\mathfrak{m})\cong\C$.
\end{theorem}

\begin{proof}
Complexifying, $\mathfrak{m}_\C=\mathbf{16}_{-3}\oplus\overline{\mathbf{16}}_{+3}$
with non-isomorphic summands exchanged by conjugation. A real submodule would
be a complex submodule invariant under conjugation; Schur's lemma on each
summand forbids nontrivial proper submodules.
\end{proof}

\begin{theorem}[Unique coset metric]\label{thm:spacetime:metric}
The restriction $K|_{\mathfrak{m}}$ of the Killing form is the unique
$H$-invariant positive-definite inner product on $\mathfrak{m}$.
\end{theorem}

\begin{proof}
$H$-invariant bilinear forms correspond to $H$-invariant maps
$\mathfrak{m}\otimes\mathfrak{m}\to\R$. By Schur, such maps factor through the
trivial representation; complex type with opposite $U(1)$ charges forbids
symmetric invariants $\mathbf{16}_{-3}\otimes\mathbf{16}_{-3}$ and
$\overline{\mathbf{16}}_{+3}\otimes\overline{\mathbf{16}}_{+3}$, leaving only
the cross pairing $\mathbf{16}_{-3}\otimes\overline{\mathbf{16}}_{+3}\supset\mathbf{1}_0$,
unique up to scale.
\end{proof}

\begin{remark}
$K|_{\mathfrak{m}}$ is \textbf{positive-definite}, not Lorentzian. Signature
$(-,+,+,+)$ is carried by the soldering form, not by the Killing restriction.
\end{remark}

%----------------------------------------------------------------------%
\subsection{Postulate P1: soldering}
\label{sec:spacetime:soldering}
%----------------------------------------------------------------------%

The $2\times2$ Hermitian-octonionic block embeds as $H_2(\Oct)\cong\R^{1,9}$
with $\det$ the Minkowski norm; $\mathrm{SL}(2,\Oct)\cong\Spin(1,9)$.

\begin{postulate}[Soldering and dimension, P1]\label{post:spacetime:P1}
There exists a soldering map $\sigma$ identifying four clock-field directions
$e^a_\mu$ ($a=0,1,2,3$) with a four-dimensional subspace of $H_2(\Oct)$ carrying
$\eta_{ab}=\mathrm{diag}(-1,+1,+1,+1)$. The world manifold has $d=4$. The
vierbein emerges from the Maurer--Cartan $\mathfrak{m}$-component
(Section~\ref{sec:gauge-gravity}).
\end{postulate}

\begin{proposition}[Emergent vierbein]\label{prop:spacetime:vierbein}
For a coset section $g(x)\in\E{6}$ with $g^{-1}\partial_\mu g=A_\mu+e_\mu$,
the $\mathfrak{m}$-component defines a composite vierbein. Under P1,
$\eta_{\mu\nu}=e^a_\mu e^b_\nu\eta_{ab}$ has signature $(-,+,+,+)$.
\end{proposition}

%----------------------------------------------------------------------%
\subsection{Emergent time}
\label{sec:spacetime:time}
%----------------------------------------------------------------------%

\begin{proposition}[Page--Wootters clock]\label{prop:spacetime:time}
In the semiclassical regime of induced gravity
(Theorem~\ref{thm:gauge-gravity:planck}), the Hamiltonian constraint
$H_{\mathrm{phys}}|\psi\rangle=0$ holds on physical states. The soldered timelike
clock $\orderparam^0$ defines internal time; conditional evolution
$U(\tau)=e^{-iH_\tau\tau}$ on physical states recovers standard quantum dynamics.
\end{proposition}

\begin{remark}[OS reconstruction]
Euclidean correlators follow from \eqref{eq:dynamics:action}. Lorentzian unitary
dynamics requires Osterwalder--Schrader reconstruction along $\orderparam^0$.
\textbf{Conjecture:} the $\sigma$-model satisfies reflection positivity.
\end{remark}% --- end 05-spacetime.tex ---
% --- 06-gauge-gravity.tex ---
%======================================================================%
\section{Gauge Sector and Induced Gravitation}
\label{sec:gauge-gravity}
%======================================================================%

%----------------------------------------------------------------------%
\subsection{Composite connection}
\label{sec:gauge-gravity:gauge}
%----------------------------------------------------------------------%

The Lie algebra splits $H$-equivariantly as
$\mathfrak{e}_6=\mathfrak{h}\oplus\mathfrak{m}$ with
$\mathfrak{h}=\mathfrak{so}_{10}\oplus\mathfrak{u}_1$ and $\mathfrak{m}$ as in
\eqref{eq:foundations:adjoint}.

\begin{theorem}[Composite gauge connection]\label{thm:gauge-gravity:composite}
For $g(x)\in\E{6}$, write $g^{-1}\partial_\mu g=\mathcal{A}_\mu+e_\mu$ with
$\mathcal{A}_\mu\in\mathfrak{h}$, $e_\mu\in\mathfrak{m}$. Under
$g\mapsto gh^{-1}$, $h\in H$,
\[
  \mathcal{A}_\mu\mapsto h\mathcal{A}_\mu h^{-1}-(\partial_\mu h)h^{-1},
\]
and the curvature $\mathcal{F}_{\mu\nu}=\partial_{[\mu}\mathcal{A}_{\nu]}+
[\mathcal{A}_\mu,\mathcal{A}_\nu]$ transforms as an $H$ field strength. All
$46$ generators of $\mathfrak{h}$ remain massless because $H$ is unbroken.
\end{theorem}

\begin{proof}
Standard transformation of the canonical connection on the reductive homogeneous
bundle $\E{6}\to\mathcal{M}$.
\end{proof}

\begin{proposition}[Induced Yang--Mills]\label{prop:gauge-gravity:ym}
At one loop, integrating out massive coset modes generates
$S_{\mathrm{YM}}=\frac{1}{4\hat g^2}\int\Tr(\mathcal{F}_{\mu\nu}\mathcal{F}^{\mu\nu})$
with $\hat g^{-2}\propto f^2\mu^2/(16\pi^2)$, $[f]=\mathrm{mass}$.
\end{proposition}

%----------------------------------------------------------------------%
\subsection{Anomaly cancellation}
\label{sec:gauge-gravity:anomaly}
%----------------------------------------------------------------------%

Anomalies are computed with Dynkin indices $A_{\mathbf{16}}=A_{\mathbf{10}}=1$,
$A_{\mathbf{1}}=0$ and $U(1)$ charges from \eqref{eq:foundations:27branch}.

\begin{theorem}[Green--Schwarz cancellation]\label{thm:gauge-gravity:gs}
The cubic anomaly $\mathcal{A}_{U(1)^3}$ vanishes:
$16\cdot1^3+10\cdot(-2)^3+1\cdot4^3=0$.
The mixed anomaly $\mathcal{A}_{U(1)\cdot\Spin(10)^2}$ has coefficient $-1$:
$1\cdot1+(-2)\cdot1+4\cdot0=-1$,
and is cancelled by a Green--Schwarz term $\propto\mathrm{tr}(F\wedge F)\,B$
with axion $B$ from the $\mathbf{1}_{+4}$ sector.
\end{theorem}

\begin{proof}
Direct arithmetic on charges and indices; see~\cite{GreenSchwarz1984} for the
mechanism.
\end{proof}

%----------------------------------------------------------------------%
\subsection{Induced gravity}
\label{sec:gauge-gravity:gravity}
%----------------------------------------------------------------------%

\begin{theorem}[Sakharov induced gravity]\label{thm:gauge-gravity:planck}
In a heat-kernel scheme, only massive modes contribute to the $a_2$ coefficient.
With $28$ physical massive coset scalars ($s=+1$, $m^2=\mu^2$ canonically),
\begin{equation}\label{eq:gauge-gravity:planck}
  \MPl^2
  =\frac{1}{16\pi^2}\cdot\frac{1}{6}\sum_{\mathrm{massive}} s_i m_i^2
  \log\frac{\Lambda^2}{m_i^2}
  \approx\frac{28}{6}\cdot\frac{\mu^2}{16\pi^2}
  \log\frac{\Lambda^2}{\mu^2}>0.
\end{equation}
The four soldering directions are protected by diffeomorphism and local Lorentz
Ward identities, not by $\Hess V$.
\end{theorem}

\begin{remark}[Cosmological constant]
The $a_0$ term yields an uncancelled $\Lambda^4$ vacuum energy. SJET does not
solve the cosmological constant problem; $\Lambda_{\mathrm{cc}}$ is a tuned
relevant coupling.
\end{remark}

%----------------------------------------------------------------------%
\subsection{Renormalization group}
\label{sec:gauge-gravity:rg}
%----------------------------------------------------------------------%

\begin{theorem}[Gauge $\beta$-function]\label{thm:gauge-gravity:beta-g}
With canonical dimensionless coupling $\hat g^2$,
\begin{equation}\label{eq:gauge-gravity:beta}
  \beta_{\hat g^2}=2\hat g^2+\frac{b_0}{48\pi^2}\hat g^4,
  \qquad
  b_0=\tfrac{11}{3}C_A-\tfrac{4}{3}T\,n_F-\tfrac{1}{6}T\,n_S.
\end{equation}
For three chiral $\mathbf{16}$s ($n_F=3$) and $n_S=28$ coset scalars:
\begin{equation}\label{eq:gauge-gravity:b0}
  b_0(3\,\mathrm{gen})
  =\tfrac{11}{3}(8)-\tfrac{4}{3}(2)(3)-\tfrac{1}{6}(2)(28)
  =\frac{88-24-28}{3}=12.
\end{equation}
The one-generation anchor ($n_F=1$, $n_S=2$) gives $b_0(1\,\mathrm{gen})=26$.
The gauge sector is asymptotically free ($\hat g^2_\star=0$).
\end{theorem}

\begin{conjecture}[Gravitational UV fixed point]\label{conj:gauge-gravity:kappa}
A truncated RG gives $\beta_\kappa=2\kappa-(5/48\pi^2)\kappa^2$ with
UV-attractive $\kappa_\star=96\pi^2/5$. This one-loop result is indicative,
not rigorous~\cite{Reuter1998}.
\end{conjecture}% --- end 06-gauge-gravity.tex ---
% --- 07-standard-model.tex ---
%======================================================================%
\section{Standard Model Descent}
\label{sec:sm}
%======================================================================%

Matter and electroweak symmetry breaking are \textbf{postulated} (P2--P4).
Anomaly cancellation is a \textbf{theorem} once the $\mathbf{27}$ content is
fixed.

%----------------------------------------------------------------------%
\subsection{Postulates P2--P4}
\label{sec:sm:postulates}
%----------------------------------------------------------------------%

\begin{postulate}[Chiral matter, P2]\label{post:sm:P2}
Each family furnishes one chiral $\mathbf{16}$ of $\Spin(10)$ with charges from
\eqref{eq:foundations:27branch}. No mirror $\overline{\mathbf{16}}$ partners.
\end{postulate}

\begin{postulate}[Three families, P3]\label{post:sm:P3}
Horizontal $SU(3)_F$ acts on $\jalg\otimes\C^3$; the three factors are three
chiral $\mathbf{16}$s. Three generations are \emph{not} derived from
Axioms~I--III.
\end{postulate}

\begin{postulate}[GUT--Higgs and Yukawa, P4]\label{post:sm:P4}
Breaking proceeds via $\mathbf{45}_H$ or $\mathbf{54}_H$
($\Spin(10)\to SU(5)\times U(1)_X$), $\mathbf{126}_H$ (breaks $B-L$), and
$\mathbf{10}_H$ (electroweak). Yukawas $y_{ij}\,\mathbf{16}_i\mathbf{16}_j H$
generate fermion masses; type-I seesaw from $\mathbf{126}_H$ gives
$m_\nu\sim m_D^2/M_R$.
\end{postulate}

\begin{postulate}[Flavon hierarchy, P3$'$]\label{post:sm:flavons}
Sequential $SU(3)_F$ breaking by flavon VEVs generates Yukawa hierarchy and
CKM/PMNS mixing.
\end{postulate}

%----------------------------------------------------------------------%
\subsection{Descent chain}
\label{sec:sm:descent}
%----------------------------------------------------------------------%

\begin{theorem}[SM gauge structure under P4]\label{thm:sm:descent}
Assume P4. The breaking chain
\[
  \E{6}\xrightarrow{V}\Spin(10)\times U(1)
  \xrightarrow{\mathbf{45}_H}\mathrm{SU}(5)\times U(1)_X
  \xrightarrow{\mathbf{126}_H}\mathrm{SU}(3)_c\times\mathrm{SU(2)}_L\times U(1)_Y
  \xrightarrow{\mathbf{10}_H}\mathrm{SU}(3)_c\times U(1)_{\mathrm{em}}
\]
yields the SM gauge group with GUT normalization $\alpha_1=(5/3)\alpha_Y$.
\end{theorem}

Per family $\mathbf{16}=(q,u^c,d^c,\ell,e^c,\nu^c)$; $\mathbf{10}_H$ supplies
$v=246\,\mathrm{GeV}$.

\begin{remark}[Yukawa source term]
A bare coupling $\bar\psi\,\orderparam^a T_a\,\psi$ with $\orderparam$ in
$\{\mathbf{16},\overline{\mathbf{16}}\}$ is \emph{not} $H$-invariant and does not
generate masses at the $\orderparam_\star$ vacuum. Masses require P4 Higgs
multiplets~\cite{GeorgiGlashow1974}.
\end{remark}

%----------------------------------------------------------------------%
\subsection{Scalar spectrum}
\label{sec:sm:masses}
%----------------------------------------------------------------------%

\begin{theorem}[Tree-level Goldstones]\label{thm:sm:tree}
At $\orderparam_\star\in\mathcal{M}$, all $32$ coset directions are massless at
tree level (Corollary~\ref{cor:dynamics:goldstone}).
\end{theorem}

\begin{theorem}[Radiative degeneracy]\label{thm:sm:radiative}
The Coleman--Weinberg potential~\cite{ColemanWeinberg1973} is $H$-invariant.
Since $\mathfrak{m}$ is a single real-irreducible $H$-module
(Theorem~\ref{thm:spacetime:irrep}), Schur gives one eigenvalue on all $28$
physical massive scalars; canonical mass $m^2=\mu^2$.
\end{theorem}

\begin{remark}[Heterotic corroboration]\label{rem:sm:heterotic}
The heterotic $\mathbb{Z}_5$ quintic gives $h^1(\hat Q,V_{16})=3$~\cite{Witten1985}.
\textbf{Conjecture}, not a derivation of P3.
\end{remark}% --- end 07-standard-model.tex ---
% --- 08-verification.tex ---
%======================================================================%
\section{Verification Ledger}
\label{sec:verification}
%======================================================================%

This section records epistemic status and independent checks. The ledger is
the referee-facing summary of what is proved, assumed, or conjectured.

%----------------------------------------------------------------------%
\subsection{Status table}
\label{sec:verification:ledger}
%----------------------------------------------------------------------%

\begin{table}[ht]
\centering
\small
\renewcommand{\arraystretch}{1.15}
\begin{tabular}{@{}p{0.28\linewidth}p{0.12\linewidth}p{0.22\linewidth}p{0.18\linewidth}@{}}
\toprule
\textbf{Claim} & \textbf{Status} & \textbf{Check} & \textbf{Ref.} \\
\midrule
Stone representation & Theorem & Standard & \S\ref{sec:foundations:axiomI} \\
Hurwitz ladder & Theorem & dims $1,2,4,8$ & \ref{thm:ladder:hurwitz} \\
$1|10|16$ partition & Prop. & Voronoi masses & \ref{prop:foundations:partition} \\
Lift $\cong\jalg$ & Theorem & 27 coords & \ref{thm:foundations:stone-lift} \\
Orbit classification & Theorem & rank, $N$ & \ref{thm:foundations:orbits} \\
Potential from $\capacity$ & Theorem & limit & \ref{thm:dynamics:potential} \\
EIII vacuum & Theorem & coercivity & \ref{thm:dynamics:vacuum} \\
$\mathfrak{m}$ irreducible & Theorem & Schur & \ref{thm:spacetime:irrep} \\
Unique coset metric & Theorem & charge pairing & \ref{thm:spacetime:metric} \\
Composite gauge & Theorem & Maurer--Cartan & \ref{thm:gauge-gravity:composite} \\
$U(1)^3$ anomaly $=0$ & Theorem & $16-80+64$ & \ref{thm:gauge-gravity:gs} \\
Mixed GS anomaly & Theorem & coeff.\ $-1$ & \ref{thm:gauge-gravity:gs} \\
$\MPl^2>0$ & Theorem & 28 scalars & \eqref{eq:gauge-gravity:planck} \\
$b_0(3)=12$ & Theorem & \eqref{eq:gauge-gravity:b0} & \ref{thm:gauge-gravity:beta-g} \\
Asymptotic freedom & Theorem & $b_0>0$ & \ref{thm:gauge-gravity:beta-g} \\
CW mass degeneracy & Theorem & Schur & \ref{thm:sm:radiative} \\
\midrule
Soldering $(1,3)$, $d=4$ & Postulate & P1 & \ref{post:spacetime:P1} \\
Chiral $\mathbf{16}$, 3 families & Postulate & P2--P3 & \S\ref{sec:sm:postulates} \\
GUT--Higgs/Yukawa & Postulate & P4 & \ref{post:sm:P4} \\
$\Lambda_{\mathrm{cc}}$ & Postulate & tuned & \S\ref{sec:gauge-gravity:gravity} \\
\midrule
SM gauge (full masses) & Under P4 & descent & \ref{thm:sm:descent} \\
$\kappa_\star=96\pi^2/5$ & Conjecture & 1-loop RG & \ref{conj:gauge-gravity:kappa} \\
OS positivity & Conjecture & $\sigma$-model & \S\ref{sec:spacetime:time} \\
Heterotic $h^1=3$ & Conjecture & string & \ref{rem:sm:heterotic} \\
\bottomrule
\end{tabular}
\caption{Epistemic ledger for SJET.}
\label{tab:verification:ledger}
\end{table}

%----------------------------------------------------------------------%
\subsection{Algebraic checks}
\label{sec:verification:algebra}
%----------------------------------------------------------------------%

\paragraph{Anomalies.}
Cubic: $16-80+64=0$. Mixed: $1-2+0=-1$ (GS).

\paragraph{$\beta$-function.}
$b_0=(88-24-28)/3=12$; one generation: $(88-8-4)/3=26$.

\paragraph{Dimensions.}
$45+1+16+16=78$; $1+10+16=27$; $16+16=32$; $28+4=32$.

%----------------------------------------------------------------------%
\subsection{Falsifiable predictions}
\label{sec:verification:predictions}
%----------------------------------------------------------------------%

\begin{enumerate}[label=\textbf{F\arabic*.},leftmargin=*]
  \item \textbf{Proton decay.} $\tau(p\to e^+\pi^0)\sim\MGUT^4/(\alpha_{\mathrm{GUT}}^2 m_p^5)$;
    Super-K/Hyper-K constrain $\MGUT\gtrsim10^{34}\,\mathrm{GeV}$.
  \item \textbf{Neutrino mass scale.} Seesaw with $M_R\sim10^{14}\,\mathrm{GeV}$ gives
    $m_\nu\sim0.05\,\mathrm{eV}$.
  \item \textbf{No fourth generation.} A fourth $\mathbf{16}$ shifts $b_0$ and
    threatens asymptotic freedom.
  \item \textbf{Coset mass degeneracy.} One-loop CW mass $\mu$ on all $28$ scalars;
    splittings only below $\MGUT$ from $H$-breaking.
  \item \textbf{$R^2$ bounds.} If Conjecture~\ref{conj:gauge-gravity:kappa} holds,
    higher-curvature couplings are suppressed in the UV.
\end{enumerate}% --- end 08-verification.tex ---
% --- 09-conclusions.tex ---
%======================================================================%
\section{Conclusions}
\label{sec:conclusions}
%======================================================================%

Stone--Jordan Exceptional Theory is a \emph{conditional} unification program:
three axioms fix the discrete--algebraic arena and capacity dynamics; postulates
P1--P4 complete the embedding into four-dimensional quantum field theory.

\subsection{Proved from axioms (modulo standard mathematics)}

\begin{itemize}
  \item Identification of the capacity limit with the sharp-map potential
    \eqref{eq:dynamics:V}.
  \item Vacuum selection on $\mathcal{M}=\mathrm{EIII}$; origin is unstable.
  \item Irreducibility and unique metric on $\mathfrak{m}$.
  \item Composite $\Spin(10)\times U(1)$ connection and anomaly cancellation.
  \item Induced $\MPl^2>0$ from $28$ massive scalars; $b_0(3)=12$.
\end{itemize}

\subsection{Postulated for SM contact}

P1 (soldering), P2--P4 (matter, families, Higgs/Yukawa), and tuned
$\Lambda_{\mathrm{cc}}$.

\subsection{Open problems}

\begin{enumerate}[label=\textbf{O\arabic*.},leftmargin=*]
  \item Rigorous gravitational UV fixed point beyond one-loop truncation.
  \item OS reflection positivity for \eqref{eq:dynamics:action}.
  \item Derive three generations without P3.
  \item Heterotic embedding of the exceptional cascade.
  \item Cosmological constant mechanism.
\end{enumerate}

SJET is mathematically closed within its stated epistemic labels, internally
consistent on anomaly and RG arithmetic, and verifiable via
Table~\ref{tab:verification:ledger} and predictions F1--F5.% --- end 09-conclusions.tex ---
% --- 10-appendix.tex ---
%======================================================================%
\appendix
\section{Albert Algebra Identities}
\label{sec:appendix}
%======================================================================%

%----------------------------------------------------------------------%
\subsection{Matrix model}
\label{sec:appendix:matrix}
%----------------------------------------------------------------------%

Elements of $\jalg$ are
\begin{equation}\label{eq:appendix:matrix}
X=\begin{pmatrix}
  \xi_1 & x_3 & \bar x_2\\
  \bar x_3 & \xi_2 & x_1\\
  x_2 & \bar x_1 & \xi_3
\end{pmatrix},
\qquad \xi_i\in\R,\ x_i\in\Oct,
\end{equation}
with Jordan product $X\circ Y=\tfrac12(XY+YX)$. The trace inner product is
$\langle X,Y\rangle=\Tr(X\circ Y)$ and $Q(X)=\Tr(X^2)$.

%----------------------------------------------------------------------%
\subsection{Sharp map and rank}
\label{sec:appendix:sharp}
%----------------------------------------------------------------------%

The cubic norm is $N(X)=\det X$ (octonionic determinant). The sharp map is
\begin{equation}\label{eq:appendix:sharp}
X^{\#}=\adj(X),\qquad X^{\#}\circ X=N(X)\,E,\qquad (X^{\#})^{\#}=N(X)\,X,
\end{equation}
with $E=\mathrm{diag}(1,1,1)$. In diagonal form
$X=\mathrm{diag}(\lambda_1,\lambda_2,\lambda_3)$,
$X^{\#}=\mathrm{diag}(\lambda_2\lambda_3,\lambda_3\lambda_1,\lambda_1\lambda_2)$.
Thus $\rank X\leq1\iff X^{\#}=0$.

%----------------------------------------------------------------------%
\subsection{Freudenthal identity}
\label{sec:appendix:freudenthal}
%----------------------------------------------------------------------%

For the quadratic adjoint $X\mapsto X^{\#}$,
\begin{equation}
  (X\circ Y)^{\#}=(N(X)Y+N(Y)X-X^{\#}\circ Y^{\#}),
\end{equation}
which characterizes $\jalg$ among cubic vector spaces. Any $\E{6}$-invariant
quartic vanishing on $\{X^{\#}=0\}$ is proportional to $\langle X^{\#},X^{\#}\rangle$.% --- end 10-appendix.tex ---
%======================================================================%
\section*{Acknowledgments}
\addcontentsline{toc}{section}{Acknowledgments}
%======================================================================%

The author thanks prior UCFT/BAMLP development for the validated spine
(archived in \texttt{old/}), especially the corrected $\beta$-function counts
and anomaly arithmetic.

%======================================================================%
\section*{References}
\addcontentsline{toc}{section}{References}
%======================================================================%

\begin{thebibliography}{99}

\bibitem{Jacobson1968}
N.~Jacobson, \emph{Structure and Representations of Jordan Algebras},
Amer.\ Math.\ Soc., Providence, 1968.

\bibitem{Freudenthal1954}
H.~Freudenthal, ``Beziehungen der $\mathfrak{E}_7$ und $\mathfrak{E}_8$
zur Oktavenebene,'' \emph{Indag.\ Math.} \textbf{16} (1954), 218--230.

\bibitem{Hurwitz1898}
A.~Hurwitz, ``\"Uber die Komposition der quadratischen Formen von
unendlich vielen Variablen,'' \emph{Math.\ Ann.} \textbf{54} (1898), 1--26.

\bibitem{Stone1936}
M.~H.~Stone, ``The theory of representations for Boolean algebras,''
\emph{Trans.\ Amer.\ Math.\ Soc.} \textbf{40} (1936), 37--111.

\bibitem{Sakharov1967}
A.~D.~Sakharov, ``Vacuum quantum fluctuations in curved space and the
theory of gravitation,'' \emph{Dokl.\ Akad.\ Nauk SSSR} \textbf{177} (1967),
70--71.

\bibitem{GreenSchwarz1984}
M.~B.~Green and J.~H.~Schwarz, ``Anomaly cancellations in supersymmetric
$D=10$ gauge theory and superstring theory,''
\emph{Phys.\ Lett.\ B} \textbf{149} (1984), 117--122.

\bibitem{GeorgiGlashow1974}
H.~Georgi and S.~L.~Glashow, ``Unity of all elementary-particle forces,''
\emph{Phys.\ Rev.\ Lett.} \textbf{32} (1974), 438--441.

\bibitem{PageWootters1983}
D.~N.~Page and W.~K.~Wootters, ``Evolution without evolution: an explicit
example,'' \emph{Phys.\ Rev.\ D} \textbf{27} (1983), 2885--2892.

\bibitem{Reuter1998}
M.~Reuter, ``Nonperturbative evolution equation for quantum gravity,''
\emph{Phys.\ Rev.\ D} \textbf{57} (1998), 971--985.

\bibitem{ColemanWeinberg1973}
S.~Coleman and E.~Weinberg, ``Radiative corrections as the origin of
spontaneous symmetry breaking,'' \emph{Phys.\ Rev.\ D} \textbf{7} (1973),
1888--1910.

\bibitem{Witten1985}
E.~Witten, ``Symmetry breaking patterns in superstring models,''
\emph{Nucl.\ Phys.\ B} \textbf{258} (1985), 75--100.

\end{thebibliography}\end{document}
